\section{Conclusion}
\label{sec:conclusion}

We presented ConjLean and REFUTE, complementary systems for automated
mathematical conjecture generation, Lean~4 formalization, and multi-agent
counterexample refutation.
On a 92-conjecture benchmark, REFUTE with Qwen2.5-Math-7B achieves perfect
precision ($1.000$) and recall $0.292$, with all refutations attributable to
deterministic boundary sweep.
LLM-dependent strategies (symbolic perturbation, analogical transfer) do not
yet contribute at 7B scale.

Scaling from Qwen2.5-Math-7B to Qwen2.5-32B-Instruct (NF4 int4) produces no
change in recall, precision, or the set of refuted conjectures.
This result strongly suggests the bottleneck is not LLM capacity for candidate
\emph{generation} but for candidate \emph{validity}: the LLM-dependent
strategies (SYMBOLIC\_PERTURBATION, ANALOGICAL) produce candidates that are
structurally plausible but fail SymPy verification at both scales.
Fine-tuning on verified counterexample triples, rather than general
instruction-following data, is the indicated next step.

Key lessons from E1 and E2:
(1) zero false positives are achievable through exact SymPy verification, a
property that scales model scale cannot erode;
(2) the recall bottleneck lies squarely in LLM-generated candidate quality for
non-trivial strategies;
(3) a correct Strategist early-stop policy matters for experimental integrity
even when it does not change aggregate recall.

\paragraph{Future work.}
Domain-specific LoRA fine-tuning of the Refuter on verified counterexample
triples (the \texttt{make finetune-refuter} pipeline) is the most direct path
to improving SYMBOLIC\_PERTURBATION yield.
Richer C-Agent prompting that explicitly checks whether refined statements
still fail at $n=1$ would address the refinement-loop failure mode identified
in \S\ref{sec:analysis}.
Extending the benchmark to cover algebraic geometry and combinatorial
identities would stress-test ANALOGICAL more meaningfully.
